% Self-contained report fragment. Requires amsmath, amssymb, amsthm and hyperref.
% The parent document supplies theorem/lemma environments if desired; none are required here.
\section{The PA194QN6 construction and its consistency calibration}
\label{sec:qn6}

PA194QN6 combines a different axiom-emission order, coefficient-six formula
coding, and a no-push modus-ponens controller. Its Shift command retains the
sum-one test. This is a new verified composition of existing ideas. It is
not a claim to a new logical calculus, a previously unknown premise-exchange
law, a state-count reduction, or a proof-assistant endpoint.
The literal table has 194 working states and 388 rows. Its start is 0 and
its separate halt label is 194. The tape is binary, two-way infinite and
initially all zero; the head begins at coordinate zero. All moves are one
cell left or right. The native presentation starts at 74 and halts at 254.

The native SHA-256 is
\begin{quote}\small\texttt{372d68753d206e5396312951a6be777179757598ad1384784db45764efd61fbc}.
\end{quote}
The dense SHA-256 is
\begin{quote}\small\texttt{72a35e528c1061f61e211d0bcf3a954859d2dbbcfd6b9f7846e4f970bb04e4e9}.
\end{quote}
The desired ordinary mathematical theorem is
\[
 \mathrm{PA}\ \vdash\quad
 \operatorname{Halt}(\mathrm{PA194QN6},\mathrm{blank})
 \ \longleftrightarrow\ 
 \exists e\,\operatorname{Prf}_{\mathrm{PA}}(e,\ulcorner0=1\urcorner).
 \tag{QN1}
\]
Negating both sides gives the requested equivalence between nonhalting and
$\operatorname{Con}(\mathrm{PA})$. No assumption of consistency is used to
prove the equivalence. Actual nonhalting follows from consistency.

\subsection{The exact numerical calculus}

Put $p(a,b)=(a+b)(a+b+1)/2+b$ and define
\[
 I(a,b)=6p(a,b),\quad M(a,b)=6p(a,b)+1,\quad
 U(a,B)=6p(B,a)+2,\quad E(a,b)=6p(a,b)+3.
\]
Implication in the following displays means $I$. Zero is the truth code,
one is $M(0,0)$, and $N(B)=I(B,1)$ is negation. Define
$B\wedge C=N(I(B,N(C)))$ and
$\exists x_aB=N(U(a,N(B)))$.
The axiom schemas, for arbitrary natural parameters, are
\begin{align*}
 L_1(D,H)&=(\exists x_D\,U(D,H))\to H,\\
 L_2(D,A,H)&=U(D,U(A,H))\to U(A,U(D,H)),\\
 L_3(D,A,H)&=M(H,D)\to(E(H,A)\to M(A,D)),\\
 L_4(D,A,H)&=E(H,D)\to(M(A,H)\to M(A,D)),\\
 L_5(D,H)&=\exists x_D E(H,D),\\
 L_6^{\mathrm Q}(D,A,H)&=U(D,A\to H)\to(U(D,A)\to U(D,H)),\\
 \Gamma(B)&=B\to U(B,B),\\
 \rho(D,A,H,u)&=((H\to u)\to D)\to((D\to H)\to(A\to H)),
 \qquad u\le D.
\end{align*}
For the final guarded schema define
\begin{align*}
 V(a,b,c)&=M(a,b)\to U(c,M(a,b)),\\
 P(a,b,c)&=c\to M(c,b),\\
 S(a,b,c)&=U\bigl(b,(M(a,b)\to P(a,b,c))\wedge(M(c,b)\to c)\bigr),\\
 Q(a,b,c)&=U(c,S(a,b,c)\to N(P(a,b,c))),\\
 G(a,b,c)&=N(V(a,b,c)\to Q(a,b,c)).
\end{align*}
The schema $G(a,b,c)$ is available when $c\ne a,b$; a failed guard emits
zero. Derivability of this failure value follows from
$\rho(0,0,0,0)=0$. Inference rules are
modus ponens and theorem generalization $B\mapsto U(i,B)$. Numerical
formulas are allowed to have free variable labels.

The only axiom replacement from the inherited coefficient-six calculus is
\[
 L_6= P_0\to(Q_0\to R_0)
 \quad\longmapsto\quad
 L_6^{\mathrm Q}=Q_0\to(P_0\to R_0),
\]
where $P_0=U(D,A)$, $Q_0=U(D,I(A,H))$, and $R_0=U(D,H)$.
The unchanged dominated implicational basis derives the exchange formula
\[
 (a\to(b\to c))\to(b\to(a\to c)).                 \tag{QN2}
\]
This derivation uses only the restricted propositional schema and modus
ponens, together with $a,b\le I(a,b)$. It uses neither quantifier
head. Thus there is no circular appeal to the axiom being replaced.
A bundled 61-line certificate gives a separate $K,S,\mathrm{MP}$ derivation;
$K$ and $S$ are themselves derived from the unchanged restricted basis.
Instantiating (QN2) in both orders gives primitive-recursive,
conclusion-preserving proof translations between the two numerical
calculi. PA verifies them by induction on coded proofs.
The alias $I(0,0)=0$ is harmless to numerical modus ponens: all steps
require the actual constructor equality, not injective free syntax.

\subsection{Why coefficient six retains the PA calibration}

The inherited coefficient-six calibration is retained in
\texttt{references/radix\_calibration.tex}, together with its coefficient-four
and separate-diagonal-head dependencies. Two points in that proof are
necessary here.

First, recoding formulas alone is insufficient because a formula's code
also appears as a variable label in $\Gamma$. The hereditary map
\[
 F(0)=0,\qquad
 F(4p(a,b)+r)=6p(F(a),F(b))+r\quad(n>0, 0\le r<4)
\]
recodes both formulas and labels. The inverse-on-image recursion proves
injectivity. Constructor identities, bound-variable binding, diagonal
heads, and the guards $c\ne a,b$ are then preserved. $F(1)=1$.
The map is not assumed to preserve numerical order. Instead the fixed
restricted-implication proof template derives the unrestricted $\rho$
schema, using only reflexivity and domination by an implication code.
It replaces each recoded restricted axiom with a finite derivation whose
own numerical side conditions are valid at coefficient six.

Second, all natural codes must be interpreted. For $n>0$, decode
$n=6p(a,b)+r$. Interpret residues 0, 1, 2, 3 as implication, finite binary
membership, universal quantification, and equality respectively. Interpret
residues 4 and 5 as closed truth, and zero separately as truth. Recursion
terminates because constructor components are strictly smaller than the
positive code. Every source-variable label occurring in the decoded formula $B$ is below
$B$. Auxiliary bound variables in the bit predicate are renamed freshly;
in particular $x_B$ is not free in the arithmetic translation of $B$.
This proves the freshness needed by $\Gamma(B)$ even for unused residues.
The exceptional constructor identity $I(0,0)=0$ uses a fixed proof of
truth equivalent to truth implying truth.

Each guarded separation instance has a PA proof supplied by finite binary
comprehension. This supplier inserts the translated formula into a fixed
induction template; it does not decide its truth. The coincident-label
case $a=b$ uses the singleton-point matrix, rather than incorrectly
holding that variable fixed under its own binder. Annotating a numerical
derivation therefore produces a PA proof of its arithmetic translation.
Code one translates to the false bit-membership assertion
$\operatorname{Bit}(x_0,x_0)$. In the other direction the inherited
PA-to-coefficient-four refutation compiler, followed by hereditary recoding
and the preceding exchange transformation, supplies a coefficient-six
Q-calculus refutation. These are syntactic transformations verified by PA,
not applications of a truth predicate or reflection principle.

\subsection{Exact command semantics}

Write $[a;w]=p(a,w)$ for the stack constructor and use total Cantor
projections $\operatorname{hd},\operatorname{tl}$. In particular
$[0;0]=0$ and both projections of zero are zero. A store is $(t,w,d,h)$:
visible theorem, stack, offset, cached theorem. The no-push MP command is
\[
 \mathsf{MP}_x(t,w,d,h)=
 \begin{cases}
 (b,w,b,0),& I(t,b)=h,\\
 (0,w,b,0),& I(t,b)\ne h,
 \end{cases}
 \qquad b=d+x.                                      \tag{QN3}
\]
Its result and offset both retain $b$ on success; the cache becomes zero.
The result is not returned in the cache. Shift is
\[
 \mathsf S(t,w,d,h)=
 \begin{cases}
 \mathsf{halt},&t+h=1,\\
 (h,\operatorname{tl}(w),d,\operatorname{hd}(w)),&t+h\ne1.
 \end{cases}                                        \tag{QN4}
\]
The stack effect differs from the supplied Q6 machine even on a valid
proof store. Since $U(0,0)=2$ and $[2;0]=3$, applying
$\mathsf{MP}_0$ to $(0,3,0,0)$ succeeds: QN6 returns $(0,3,0,0)$,
whereas Q6 inserts its dummy zero and returns $(0,9,0,0)$ because
$[0;3]=9$. At the native-table level the change from Q6 is the single
row $(164,1)$, whose next control changes from 98 to 114 while its
write-one and right-move actions are retained. This skips a real stack
operation. Relative to N6, the 83 changed source instructions instead
implement the exchanged quantifier axiom. No supplied machine in the
recognized corpus combines both behaviors with these constructors.

For $D=d+x$, $H=h+y$, the axiom command has effect
\[
 \mathsf{AX}_{A,x,y}(t,w,d,h)
  =(g(H,D,A),[q_1,\ldots,q_9,t;w],0,0),               \tag{QN5}
\]
where $g$ is the guarded $G$, and
\[
 (q_1,\ldots,q_9)=
 (\Gamma(A),L_6^{\mathrm Q},L_5,L_4,L_3,L_2,L_1,
       \rho(D,A,H,d),U(d,h)).
\]
There are nine generated stack entries and the saved old top. Every
$q_j$ except $q_9$, and the gate, is divisible by six. The last entry
$q_9\equiv2\pmod6$ is at least two. No generated entry is code one.
These are arithmetic facts for all incoming natural stores, independent
of consistency. The Q change preserves the last two stack entries and
their residues, which are the exact properties used by the command compiler.

Let $A_0=\mathsf{AX}_{0,0,0}$ and define
\[
 \mathsf{Push}=A_0;\mathsf S^9;\mathsf{MP}_0,
 \qquad \mathsf{Pop}=A_0;\mathsf S^{11},
 \qquad \mathsf{Discard}=\mathsf{MP}_0.
\]
Directly from (QN3)--(QN5), for all displayed natural values,
\begin{align*}
 \mathsf{Push}(t,w,d,h)&=(0,[t;w],0,0),\\
 \mathsf{Pop}(t,[a;w],d,h)&=(t,w,0,a),\\
 \mathsf{Discard}(t,w,0,h)&=(0,w,0,0),\\
 \mathsf{MP}_i(0,w,0,0)&=(0,w,i,0).
\end{align*}
The first ten Shifts after AX end at $(U(d,h),w,0,t)$; the eleventh
has sum $U(d,h)+t\ge2$. Hence Pop is total even when $t=1$.
Discard is one command because when $d=x=0$ both MP outcomes coincide.
The offset contract follows from $I(0,i)>0$ for $i>0$ and is also valid
at $i=0$. These equations use the literal Cantor identities, including
zero-stack collapse; they do not identify numerical stacks with free lists.

From $(t,w,0,0)$ a selected head $q_j$, $1\le j\le7$, is obtained by
\[
 \mathsf{Discard};\mathsf{AX}_{A,D,H};
       \mathsf S^{j+1};\mathsf{Pop}^{9-j}.
\]
The dominated head at $(a,b,c,u)$, $u\le a$, is selected by
\[
 \mathsf{Discard};\mathsf{MP}_u;
       \mathsf{AX}_{b,a-u,c};\mathsf S^9;\mathsf{Pop}.
\]
The gate uses
$\mathsf{Discard};\mathsf{AX}_{A,D,H};\mathsf{Pop}^{10}$.
All return the desired formula in the top, preserve the framed stack,
and clear offset and cache. Generalization is compiled by
\[
 \mathsf{Push};\mathsf{MP}_i;\mathsf S;
      \mathsf{AX}_{i,0,0};\mathsf S^{10},
\]
which takes $(a,w,0,0)$ to $(U(i,a),w,0,0)$.

For a modus-ponens inference from $a$ and $I(a,b)$, recursively compiled
premises are combined as
\[
 \mathsf{Discard};\mathsf{Push};C(I(a,b));\mathsf{Push};
 C(a);\mathsf{Push};\mathsf S;\mathsf S;
       \mathsf{MP}_b;\mathsf{Pop}.                    \tag{QN6}
\]
Before the two Shifts the store is
$(0,[a,I(a,b),0;w],0,0)$. They either halt, if the antecedent is already
one, or give $(a,[0;w],0,I(a,b))$. MP gives $(b,[0;w],b,0)$;
the final total Pop gives $(b,w,0,0)$. The explicitly reserved zero is
consumed exactly once. In particular, zero-pair collapse and the
$b=1$ case satisfy the same equations.
Induction on a derivation proves that $C$ either halts or returns its
conclusion with arbitrary framed stack preserved and both auxiliary proof
registers cleared. Prefixing $A_0$ first allows arbitrary incoming offsets
and caches. Thus a refutation gives the word $A_0;C(e);\mathsf S$ that
halts from every incoming store.

\subsection{Scheduling, soundness, and physical refinement}

The scheduler carries its store between trials. Its argument therefore
uses the all-store compiler, not equality of the stores of different
machines after the same candidate counter. Positive continuation codes
are $E(\epsilon)=0$ and
$E(c::w)=1+\operatorname{enc}(c,E(w))$. Fresh candidate codes omit this
outer successor. In particular
\[
 \operatorname{enc}(\mathsf{AX}_{A,x,y},z)
       =p(A+2,p(y,p(x,z))).
\]
For the universal halting word $W=A_0;C(e);\mathsf S$, the prefixed
words $\mathsf{AX}_{k,0,0};W$ have fresh codes
$p(k+2,p(0,p(0,E(W))))$, a strictly increasing unbounded family.
Fresh candidates are enumerated from two; positive residuals strictly
decrease; every constituent register loop terminates. Therefore one of
these candidates is reached unless the machine has halted earlier.

For soundness, annotate the visible theorem, cached theorem, and every
entry in the canonical finite Cantor expansion of $w$ by numerical proofs
or PA proofs of their translations. AX supplies valid heads; MP composes
the two annotated premises or outputs truth; Shift transports annotations.
The test $t+h=1$ implies that one annotated theorem is code one, hence
supplies a PA refutation. Arbitrary arithmetic offsets need no annotation.
This proves both source/refutation reductions by syntactic induction in PA.

The literal implementation is checked separately. The 202-control raw
machine has 135 fixed controls and 67 AX readers. Eight simultaneous
history-sensitive folds leave 194 controls. Thirty-two complete two-cell
contracts verify the folds on the final simultaneous table, with arbitrary
untouched exterior tape. All 4,877 primitive source outcomes are checked
for arbitrary natural register values and both decrement branches.
The independent literal-tape checker rederives four register-list lemmas
from actual transition rows and assumes no arithmetic-box oracle.
Its frame permits an arbitrary nonempty finite list of natural spare values
and an arbitrary affine header gap. Exact blank simulation reaches the
entire clean source frame after 58,320 steps, in native state 84, head
$-2$, logical registers $(2,0,\ldots,0)$, and 26 spare blocks of value one.

These checks give positive finite physical simulations of source steps.
Determinism and positivity give the converse halting direction: a physical
first halt cannot precede an arbitrarily long positive simulation of a
nonhalting source. The elementary finite row calculations, unary-loop
inductions, and coded trace inductions are arithmetizable in PA.
Composing them with the syntactic source/refutation reductions proves (QN1).
The Python artifacts support this ordinary mathematical proof; they do not
constitute a new kernel-checked Lean or Rocq theorem, nor discharge the
repository's separate machine-specific iPRA goals.
